% Requires amsmath,amssymb,amsthm; theorem,lemma,corollary,remark environments.
% Bibliography keys used: reiffenstein, dlmfdirichlet.
\section{Entire functions, collisions, and zero counting}
\label{sec:entire-graph}

The finite compression theorem has an entire-function counterpart.
Its real-zero mechanism is the classical sign of a logarithmic derivative
of a real canonical product.  In matrix language this is a special case
of meromorphic Herglotz theory and its higher-order interlacing;
we do not claim that mechanism as new.\cite{reiffenstein}
We give direct proofs to make the multiplicities and all pole cancellations
in the present construction explicit.

\subsection{The precise analytic class}

Let $F_1,\ldots,F_n$ be nonzero entire functions with representations
\begin{equation}\label{eq:real-canonical-product}
 F_i(z)=c_i z^{m_i}e^{b_i z}
 \prod_{\rho\in\mathcal Z_i}
 \left(1-\frac z\rho\right)e^{z/\rho},
 \qquad
 c_i\in\mathbb R\setminus\{0\},\quad b_i\in\mathbb R,\quad m_i\geq0.
\end{equation}
Here $\mathcal Z_i$ is a possibly empty multiset of nonzero real numbers,
each zero is listed with multiplicity,
$\sum_{\rho\in\mathcal Z_i}|\rho|^{-2}<\infty$, and the product converges
locally uniformly.  Assume that each $F_i$ has at least one zero:
$m_i+|\mathcal Z_i|>0$.
These hypotheses hold, in particular, for real entire functions of order
at most one with at least one zero and with all zeros real, by Hadamard
factorization.  They also include the polynomial inputs above.
No assumption that $F_i(0)\neq0$ is made.

Write $g_i=F_i'/F_i$ away from the zeros of $F_i$.
Termwise logarithmic differentiation of
\eqref{eq:real-canonical-product} gives
\begin{equation}\label{eq:logder-canonical}
 g_i(z)=b_i+\frac{m_i}{z}
   +\sum_{\rho\in\mathcal Z_i}
       \left(\frac1{z-\rho}+\frac1\rho\right).
\end{equation}
The sum and its derivative converge uniformly on compact sets avoiding
the zeros: for large $|\rho|$ the summand in \eqref{eq:logder-canonical}
is $O(|\rho|^{-2})$, uniformly on each such compact set.

\begin{lemma}[Strict half-plane sign]\label{lem:strict-logder}
For $z=x+iy$ with $y\neq0$,
\begin{equation}\label{eq:logder-imaginary}
 \operatorname{Im}g_i(z)
 =-y\left(\frac{m_i}{|z|^2}
       +\sum_{\rho\in\mathcal Z_i}\frac1{|z-\rho|^2}\right).
\end{equation}
At real $x$ with $F_i(x)\neq0$,
\begin{equation}\label{eq:logder-decrease}
 -g_i'(x)=
 \frac{m_i}{x^2}
       +\sum_{\rho\in\mathcal Z_i}\frac1{(x-\rho)^2}>0.
\end{equation}
Terms involving $m_i$ are omitted when $m_i=0$.
\end{lemma}
\begin{proof}
The constants $b_i$ and $1/\rho$ are real, and
$\operatorname{Im}(z-\rho)^{-1}=-y/|z-\rho|^2$.
This proves \eqref{eq:logder-imaginary}.
Differentiating \eqref{eq:logder-canonical} gives
\eqref{eq:logder-decrease}.
At least one nonnegative term in each parenthesized sum is positive,
because $F_i$ has at least one zero.
\end{proof}

The analytic class matters.  Real zero location alone, without a growth
or canonical-product restriction, is insufficient: $F(z)=ze^{z^2}$ has
only a real zero, whereas $F'(z)=e^{z^2}(1+2z^2)$ has nonreal zeros.
Likewise a zero-free exponential has constant logarithmic derivative
and does not satisfy the strict inequality needed below.

\subsection{A determinant without division}

Let $A$ be an $n$-by-$n$ Hermitian matrix, constant in $z$.  For a subset
$S\subseteq[n]=\{1,\ldots,n\}$, write $A[S]$ for the corresponding
principal submatrix, with $\det A[\varnothing]=1$.  Define
\begin{align}
 \mathcal P_A(z)
 &=\det\!\left(
       \operatorname{diag}(F_1(z),\ldots,F_n(z))A+
       \operatorname{diag}(F'_1(z),\ldots,F'_n(z))
          \right) \label{eq:entire-determinant}\\
 &=\sum_{S\subseteq[n]}
       \det A[S^{\mathrm c}]
       \prod_{i\in S}F_i'(z)
       \prod_{j\notin S}F_j(z). \label{eq:principal-minor-entire}
\end{align}
The second equality follows by multilinearity in the rows of the
determinant.  It proves directly that $\mathcal P_A$ is entire.
Every principal minor of a Hermitian matrix is real, so $\mathcal P_A$
is a real entire function.
Away from the input zeros, put
\[
 M_A(z)=A+\operatorname{diag}(g_1(z),\ldots,g_n(z)).
\]
Then
\begin{equation}\label{eq:entire-meromorphic-factorization}
 \mathcal P_A(z)=\left(\prod_iF_i(z)\right)\det M_A(z).
\end{equation}

\begin{theorem}[Real zeros under Hermitian coupling]\label{thm:entire-real-zeros}
The entire function $\mathcal P_A$ is not identically zero and all of
its zeros are real.
For $A=L$, a weighted graph Laplacian, it is precisely the rooted-forest
expression \eqref{eq:forest-polynomial} with $p_i$ replaced by $F_i$.
\end{theorem}
\begin{proof}
Take $z$ in the upper half-plane.  All $F_i(z)$ are nonzero.
If $M_A(z)v=0$ for a nonzero vector $v$, Hermiticity of $A$ and
Lemma~\ref{lem:strict-logder} give
\[
 0=\operatorname{Im}\bigl(v^*M_A(z)v\bigr)
   =\sum_i|v_i|^2\operatorname{Im}g_i(z)<0,
\]
a contradiction.  Thus \eqref{eq:entire-meromorphic-factorization}
is nonzero throughout the upper half-plane.  Since $\mathcal P_A$ is
real entire, it is nonzero throughout the lower half-plane as well.
The forest assertion follows from the polynomial identity
\eqref{eq:forest-det}; its substitution is first valid off the input
zeros, and the entire expressions then agree everywhere.
\end{proof}

\subsection{An exact formula at collisions}

A product formula written with logarithmic derivatives can conceal
cancellation at a shared zero.  The following formula resolves every
such collision, including multiple and central zeros.

Fix $\lambda\in\mathbb R$, and define
\[
 I=\{i:F_i(\lambda)=0\},\qquad J=[n]\setminus I,\qquad
 \mu_i=\operatorname{ord}_{\lambda}F_i\quad(i\in I).
\]
Define the complementary Hermitian matrix
\begin{equation}\label{eq:collision-complement}
 C_\lambda=A[J]+\operatorname{diag}(g_j(\lambda):j\in J),
 \qquad k_\lambda=\dim\ker C_\lambda.
\end{equation}
If $J$ is empty, its nullity is defined to be zero.

\begin{theorem}[Exact collision order]\label{thm:collision-order}
For every real $\lambda$,
\begin{equation}\label{eq:collision-order}
 \boxed{\quad
 \operatorname{ord}_{\lambda}\mathcal P_A
    =\sum_{i\in I}(\mu_i-1)+k_\lambda .
 \quad}
\end{equation}
When $I$ is empty, this says that the multiplicity of a zero away from
the input zeros equals the nullity of $M_A(\lambda)$.
\end{theorem}
\begin{proof}
Set $t=z-\lambda$.  For $i\in I$ write
$F_i(\lambda+t)=t^{\mu_i}h_i(t)$ with $h_i(0)\neq0$.
The $I$-by-$I$ block of $M_A(\lambda+t)$ is
\[
 B(t)=t^{-1}D+E(t),\qquad
 D=\operatorname{diag}(\mu_i:i\in I),
\]
where $E$ is analytic near zero.  Hence
\[
 B(t)^{-1}=tD^{-1}+O(t^2),\qquad
 \det B(t)=t^{-|I|}\det D\,(1+O(t)).
\]
The Schur complement of this block is analytic near zero:
\begin{align*}
 S(t)&=A[J]+\operatorname{diag}(g_j(\lambda+t):j\in J)
                   -A[J,I]B(t)^{-1}A[I,J]\\
     &=C_\lambda+tD_\lambda+O(t^2),
\end{align*}
where
\begin{equation}\label{eq:collision-negative-derivative}
 D_\lambda=
   \operatorname{diag}(g_j'(\lambda):j\in J)
        -A[J,I]D^{-1}A[I,J]\prec0
\end{equation}
when $J$ is nonempty.  Indeed the first term is strictly negative
definite by \eqref{eq:logder-decrease}, and the term subtracted from it
is positive semidefinite.

We verify explicitly that $\det S(t)$ has order $k_\lambda$.
Use the orthogonal decomposition
$\mathbb C^J=\ker C_\lambda\oplus(\ker C_\lambda)^\perp$.
On the second summand $C_\lambda$ is invertible.  Taking its Schur
complement reduces the first block to
\[
 t\,D_\lambda|_{\ker C_\lambda}+O(t^2).
\]
Cross-block contributions start in order $t^2$, since the off-diagonal
blocks of $C_\lambda$ in this decomposition vanish.
The compression of $D_\lambda$ to the kernel is negative definite,
so its determinant is nonzero.  Thus
$\operatorname{ord}_0\det S=k_\lambda$.
This argument also covers a zero-dimensional kernel; the empty
determinant is one.

Finally, the factors $\prod_{i\in I}F_i$ supply order
$\sum_{i\in I}\mu_i$, $\det B(t)$ removes $|I|$, and all $F_j(\lambda)$
for $j\in J$ are nonzero.  Equation
\eqref{eq:entire-meromorphic-factorization} now gives
\eqref{eq:collision-order}.
If $I$ is empty, apply the kernel-compression argument directly to
$M_A(\lambda+t)$, whose derivative is strictly negative definite.
If $J$ is empty, the first block alone gives the asserted order.
\end{proof}

No simplicity or noncoincidence assumption is hidden in this formula.
For example, if every input has a simple zero at $\lambda$, then
$I=[n]$, $J=\varnothing$, and $\mathcal P_A(\lambda)\neq0$:
the apparent common zero is removed completely.

\subsection{A bounded discrepancy on every interval}

Let $n_+(B)$ denote the number of strictly positive eigenvalues of a
Hermitian matrix $B$, counted with multiplicity.  For a nonzero entire
function $F$, let $N_F(a,b)$ count its real zeros in $(a,b)$ with
multiplicity.

\begin{theorem}[Exact zero-count identity]\label{thm:entire-count}
Let $a<b$ be real, and suppose neither endpoint is a zero of any $F_i$
or of $\mathcal P_A$.  Then
\begin{equation}\label{eq:inertia-count}
 \boxed{\quad
 N_{\mathcal P_A}(a,b)
  =\sum_iN_{F_i}(a,b)
     +n_+(M_A(a))-n_+(M_A(b)).
 \quad}
\end{equation}
In particular,
\begin{equation}\label{eq:count-discrepancy}
 \left|N_{\mathcal P_A}(a,b)-\sum_iN_{F_i}(a,b)\right|\leq n.
\end{equation}
If the interval contains no input zero, it contains at most $n$ zeros
of $\mathcal P_A$, counted with multiplicity.
\end{theorem}
\begin{proof}
On any interval without input zeros,
\[
 M_A'(x)=\operatorname{diag}(g_i'(x))\prec0.
\]
At a zero of $\det M_A$ with kernel dimension $k$, the same compression
argument used above shows that the $k$ zero eigenvalues cross from
positive to negative as $x$ increases.  Its positive inertia therefore
drops by $k$, which is exactly the multiplicity of $\mathcal P_A$ there.

Now cross an input zero $\lambda$ with notation as in
Theorem~\ref{thm:collision-order}.  For small negative $t$, the block
$B(t)$ has $|I|$ negative eigenvalues; for small positive $t$, it has
$|I|$ positive eigenvalues.  Hermitian block elimination preserves
inertia, so the remaining eigenvalues are counted by $S(t)$.
Its $k_\lambda$ zero eigenvalues cross from positive to negative by
\eqref{eq:collision-negative-derivative}; its other eigenvalues retain
their signs.  Consequently
\[
 n_+(M_A(\lambda+0))-n_+(M_A(\lambda-0))
       =|I|-k_\lambda
       =\sum_{i\in I}\mu_i
          -\operatorname{ord}_{\lambda}\mathcal P_A.
\]
There are only finitely many zeros of these nonzero entire functions
in a compact interval.  Summing these jumps, including the ordinary
zeros between poles, telescopes to \eqref{eq:inertia-count}.
Since positive inertia lies between zero and $n$, the two remaining
claims follow.
\end{proof}

The discrepancy bound is sharp even for connected graph Laplacians.
Take $F_i(z)=z$ for every $i$.  If the nonzero eigenvalues of $L$ are
$\ell_2,\ldots,\ell_n>0$, then
\[
 \mathcal P_L(z)=\prod_{j=2}^n(1+\ell_jz).
\]
An interval sufficiently close to zero contains all $n$ input zeros
and no output zero, attaining discrepancy $-n$.

\subsection{Local motion under a positive coupling change}

A local monotonicity statement survives for entire inputs, although
the finite resistance sum in \eqref{eq:resistance-sum} is not an
unregularized sum over infinitely many zeros.

\begin{corollary}[Velocity of a simple zero]\label{cor:entire-velocity}
Let $A(t)=A_0+tB$ with $B\succeq0$ Hermitian.
Suppose $x_0$ is a simple zero of $\mathcal P_{A(t_0)}$ and
$F_i(x_0)\neq0$ for every $i$.
There is a local differentiable real zero $x(t)$ with $x(t_0)=x_0$.
For any nonzero $v\in\ker M_{A(t_0)}(x_0)$,
\begin{equation}\label{eq:entire-velocity}
 x'(t_0)=
 \frac{v^*Bv}{\displaystyle\sum_i|v_i|^2(-g_i'(x_0))}\geq0.
\end{equation}
For increasing one graph-edge conductance, the numerator is
$|v_i-v_j|^2$.
\end{corollary}
\begin{proof}
The collision formula with $I=\varnothing$ shows that the kernel is
one-dimensional.  The real implicit function theorem applies because
the zero is simple.  Choose a local nonzero kernel vector and
differentiate $M_{A(t)}(x(t))v(t)=0$ at $t_0$.
Multiplication on the left by $v^*$ removes the term involving $v'$,
and gives
$v^*Bv+x'(t_0)\sum_i|v_i|^2g_i'(x_0)=0$.
The denominator in \eqref{eq:entire-velocity} is strictly positive.
\end{proof}

\subsection{Completed Dirichlet functions as inputs}

Let $\chi$ be a primitive nonprincipal Dirichlet character of conductor
$q$, with parity $a_\chi\in\{0,1\}$ determined by
$\chi(-1)=(-1)^{a_\chi}$.  Use the standard completion
\begin{equation}\label{eq:completed-dirichlet}
 \Lambda_\chi(s)=
  (q/\pi)^{(s+a_\chi)/2}
  \Gamma\!\left(\frac{s+a_\chi}{2}\right)L(s,\chi).
\end{equation}
Its entire continuation has order one, and its zeros are the
nontrivial zeros of $L(s,\chi)$.  The primitive functional equation and
conjugation identity are
\begin{equation}\label{eq:completed-functional-equation}
 \Lambda_\chi(s)=\varepsilon_\chi
                  \Lambda_{\overline\chi}(1-s),
 \qquad |\varepsilon_\chi|=1,\qquad
 \overline{\Lambda_\chi(\overline s)}
                  =\Lambda_{\overline\chi}(s).
\end{equation}
These are the standard analytic facts about primitive completions;
the functional-equation conventions are recorded, for example, in
DLMF \S25.15.\cite{dlmfdirichlet}

Choose a square root $\eta_\chi$ of $\varepsilon_\chi$, and define
\begin{equation}\label{eq:phase-real-dirichlet}
 F_\chi(z)=\eta_\chi^{-1}
             \Lambda_\chi\!\left(\tfrac12+iz\right).
\end{equation}
This is a real entire function even when $\chi$ is complex-valued.
Indeed, for every complex $z$,
\begin{align*}
 \overline{F_\chi(\overline z)}
 &=\eta_\chi
       \Lambda_{\overline\chi}\!\left(\tfrac12-iz\right)\\
 &=\eta_\chi\varepsilon_\chi^{-1}
       \Lambda_\chi\!\left(\tfrac12+iz\right)
 =F_\chi(z).
\end{align*}
Changing the choice of square root merely changes the sign of $F_\chi$,
which does not change any of the zero sets in the construction.

Under Dirichlet GRH, every zero of $F_\chi$ is real.
Hadamard factorization therefore gives
\eqref{eq:real-canonical-product}, with real $b_i$ and $c_i$:
divide its Hadamard factorization by the real canonical product; its
value and logarithmic derivative at a real regular point show that
the remaining constant and linear exponential can be chosen real.
The zero set is nonempty.  One way to see this, without a numerical
zero computation, is that a zero-free entire function of order at most
one is an exponential of a linear polynomial, whereas on the positive
real $s$-axis Stirling's formula and $L(s,\chi)\to1$ give
$\log|\Lambda_\chi(s)|=\tfrac{s}{2}\log s+O(s)$.
The latter growth cannot be that of a linear exponential.

For the trivial primitive character use the pole-removed zeta
completion
\[
 \xi(s)=\tfrac12s(s-1)\pi^{-s/2}\Gamma(s/2)\zeta(s),
 \qquad F_\zeta(z)=\xi\!\left(\tfrac12+iz\right).
\]
The functional equation $\xi(s)=\xi(1-s)$ makes $F_\zeta$ real and even,
and RH supplies the same real-zero canonical-product hypotheses.

\begin{corollary}[Graph-coupled completed Dirichlet functions]
\label{cor:dirichlet-graphs}
Assume the Dirichlet zero-location theorem used in the upstream proof.
Choose any finite list of primitive completed functions
\eqref{eq:phase-real-dirichlet}, allowing $F_\zeta$ and repetitions,
and any Hermitian coupling matrix $A$.
Then the explicitly entire function
\eqref{eq:principal-minor-entire} has only real zeros.
Its collision multiplicities are given by \eqref{eq:collision-order},
and on every interval with regular endpoints its zero count differs
from the sum of the input zero counts by at most $n$.

For a graph Laplacian, this yields an entire rooted-forest combination
of the completed functions and their first derivatives.
After the change of variable $z=-i(s-\tfrac12)$, every zero lies on
$\operatorname{Re}s=\tfrac12$.
\end{corollary}
\begin{proof}
The preceding normalization and factorization verify all hypotheses
of Theorems~\ref{thm:entire-real-zeros},
\ref{thm:collision-order}, and \ref{thm:entire-count}.
The linear change of variable maps the real $z$-axis to the critical
line.
\end{proof}

Central zeros cause no exception: their exact multiplicities are
retained through the powers $z^{m_i}$ in
\eqref{eq:real-canonical-product}.
A conjugate pair can also be combined into one even input.
Choose $\eta_{\overline\chi}=\overline{\eta_\chi}$.
Then \eqref{eq:completed-functional-equation} gives
$F_{\overline\chi}(z)=F_\chi(-z)$, so
$F_\chi(z)F_{\overline\chi}(z)$ is a real even canonical product
satisfying the same hypotheses.

Primitivity in this application is substantive.  An imprimitive
Dirichlet $L$-function includes extra Euler factors whose zeros can lie
on $\operatorname{Re}s=0$; they need not become real zeros in the
$z$-coordinate.  The constructed graph functions are explicit entire
functions with a critical-line zero theorem.  No Euler product,
automorphic interpretation, or new instance of Artin GRH is asserted.

